Let \(d=|\mathcal O|\), let \(\Phi=\theta\circ\Pi_{\mathfrak O}\), and define
\[
X_n=\sqrt n(\widehat p-p),
\qquad
T_n=\max\{\sqrt n(\widehat L-L(p)),\sqrt n(U(p)-\widehat U)\}.
\tag{1}
\]
By \(\text{lem:finite-multinomial-clt}\), whose sampling premise is \(\text{ass:iid-finite-sampling}\),
\[
X_n\rightsquigarrow G_p,
\qquad X_n=O_{\Pr}(1).
\tag{2}
\]
Moreover, \(X_n\in V_p\) almost surely: its coordinates sum to zero, and \(p_o=0\) implies \(N_o=0\) almost surely.  Assumption (i) in the theorem supplies \(r>0\) such that
\[
\Phi(p+x)-\Phi(p)=(D_L(x),D_U(x))
\quad\text{for }x\in V_p, \|x\|_2<r.
\tag{3}
\]
The two directional coordinates are continuous, positively homogeneous, and linear on each cone of a finite fan.  Compactness of the unit sphere in \(V_p\), or simply finiteness of its slopes, therefore gives a constant \(C_D<\infty\) for which
\[
|D_L(x)|\vee|D_U(x)|\leq C_D\|x\|_2
\qquad(x\in V_p).
\tag{4}
\]
Equations (2)--(4) and positive homogeneity imply
\[
T_n=H(X_n)
\tag{5}
\]
with probability tending to one.  Since \(H\) is continuous, (2) also gives \(H(X_n)\rightsquigarrow H(G_p)\).

We compare the conditional subsampling law to an independent size-\(m\) experiment, where below \(m=m_n\).  Let \(\Pr^*\) denote probability conditional on the full sample, let \(\widehat p_b\) be the empirical law of a uniformly selected size-\(m\) subsample without replacement, and put
\[
a_n=\sqrt m(\widehat p-p)=\sqrt{m/n}\,X_n.
\tag{6}
\]
On the event that both empirical laws are in the neighborhood in (3), the centered statistic defining the conditional critical value is exactly
\[
S_{n,b}=max\left\{
D_L\bigl(\sqrt m(\widehat p_b-p)\bigr)-D_L(a_n),
D_U(a_n)-D_U\bigl(\sqrt m(\widehat p_b-p)\bigr)
\right\}.
\tag{7}
\]
This follows by applying (3) at \(\widehat p_b\) and \(\widehat p\), subtracting, multiplying by \(\sqrt m\), and using positive homogeneity.

Couple the ordered indices sampled without replacement to \(m\) independent indices sampled with replacement from \(\{1,\ldots,n\}\).  A coupling failure requires a repeated with-replacement index, so the union bound gives conditional failure probability at most
\[
\frac{m(m-1)}{2n}.
\tag{8}
\]
For each with-replacement draw, maximally couple its empirical categorical law \(\widehat p\) to the population law \(p\).  The mismatch probability for one draw is
\[
\operatorname{TV}(\widehat p,p)=\frac12\|\widehat p-p\|_1.
\]
Another union bound shows that all \(m\) categories agree except on an event of conditional probability at most
\[
\frac m2\|\widehat p-p\|_1.
\tag{9}
\]
On \(\{\|X_n\|_2\leq M\}\), (9) is bounded by \(m\sqrt d M/(2\sqrt n)\).

Let
\[
X_m=m^{-1/2}\sum_{i=1}^m(\mathbf e_{O_i}-p)
\]
be formed from an independent size-\(m\) sample with law \(p\).  On the successful coupling event, \(\sqrt m(\widehat p_b-p)=X_m\), and the comparison variable in (7) becomes
\[
K_{a_n}(X_m),
\qquad
K_a(g):=\max\{D_L(g)-D_L(a),D_U(a)-D_U(g)\}.
\tag{10}
\]

We also control failure of the local representation (3).  Each coordinate of \(\mathbf e_{O_i}-p\) is centered and lies in \([-1,1]\).  For all sufficiently small \(|u|\), Taylor's formula on this bounded interval gives a constant \(C_1\), depending only on the finite alphabet, such that
\[
\mathbb E_p\exp\{u(\mathbf1\{O_i=o\}-p_o)\}\leq\exp(C_1u^2).
\]
Independence, exponential Markov inequality, optimization over \(u\), and a union bound over the \(d\) coordinates give constants \(C_2,c_2>0\) such that
\[
\Pr_p\{\|X_m\|_2>r\sqrt m/2\}\leq C_2e^{-c_2m}.
\tag{11}
\]
For fixed \(M\) and all sufficiently large \(n\), \(\|X_n\|_2\leq M\) puts \(\widehat p\) in the radius-\(r/2\) neighborhood in (3).  Equations (8)--(11) therefore bound the conditional total-variation error between \(S_{n,b}\) and \(K_{a_n}(X_m)\), on \(\{\|X_n\|_2\leq M\}\), by
\[
q_{n,M}:=\frac{m(m-1)}{2n}+\frac{m\sqrt d M}{2\sqrt n}+C_2e^{-c_2m}.
\tag{12}
\]

We next verify that the explicitly defined multinomial approximation error vanishes.  For every \(A\in\mathcal A_p\), the finite multinomial central limit theorem gives
\[
AX_m\rightsquigarrow AG_p.
\tag{13}
\]
Delete identically zero rows of \(A\).  By condition (ii) and the construction of \(\mathcal A_p\), each remaining row is nonzero on \(V_p\).  Gaussian nondegeneracy on \(V_p\) makes each corresponding marginal of \(AG_p\) non-atomic.  The boundary of every lower orthant is contained in a finite union of coordinate hyperplanes, and hence has \(AG_p\)-probability zero.  Thus (13) gives pointwise convergence of all lower-orthant probabilities.  A finite grid on a sufficiently large compact box, coordinatewise monotonicity between grid points, and tightness outside the box upgrade this convergence to uniform convergence over lower orthants.  Since \(\mathcal A_p\) is finite,
\[
\delta_m(p)\longrightarrow0.
\tag{14}
\]

On each pair of fan cones containing \(a_n\) and \(g\), the event \(K_{a_n}(g)\leq x\) is an intersection of linear inequalities.  The normals are selected endpoint slopes, fan-facet normals, and their sign changes.  Taking the finite union over cone pairs, and using the stipulated closure properties of \(\mathcal A_p\), writes this event as a finite Boolean combination of lower orthants governed by matrices in \(\mathcal A_p\).  Inclusion--exclusion therefore bounds the difference between its probability under \(X_m\) and under \(G_p\) by \(C_F\delta_m(p)\), for a finite constant \(C_F\) depending only on the fixed fan.  In addition, (4) gives the deterministic, uniform bound
\[
|K_{a_n}(g)-H(g)|
\leq |D_L(a_n)|\vee|D_U(a_n)|
\leq C_D\|a_n\|_2.
\tag{15}
\]
On \(\{\|X_n\|_2\leq M\}\), equations (6), (12), (14), and (15), together with \(m^2/n\to0\), imply, at every continuity point \(x\) of the distribution of \(H(G_p)\),
\[
\Pr^*\{S_{n,b}\leq x\}\longrightarrow
\Pr\{H(G_p)\leq x\}
\quad\text{in probability}.
\tag{16}
\]
Because \(X_n=O_{\Pr}(1)\), we remove the restriction to bounded \(X_n\) by letting \(M\to\infty\).

Let \(F(x)=\Pr\{H(G_p)\leq x\}\).  On the interior of each full-dimensional cone of the fan refined by the endpoint-comparison and zero hyperplanes, \(H\) is one homogeneous linear functional.  Proper cone boundaries lie in finitely many hyperplanes whose normals are nonzero on \(V_p\), and therefore have Gaussian probability zero.  A nonzero linear functional of \(G_p\) has a continuous Gaussian law, whereas a zero linear functional is identically zero.  It follows that
\[
F\text{ can have an atom only at }0.
\tag{17}
\]

Put \(\beta=1-\alpha\), let \(c_n=c^{\mathrm{sub}}_{n,\beta+\eta_n}\), with the lower-quantile convention \(\inf\{x:\Pr^*(S_{n,b}\leq x)\geq\beta+\eta_n\}\), and define \(q_\beta=\inf\{x:F(x)\geq\beta\}\).  Suppose first that \(q_\beta\neq0\).  For every sufficiently small \(\varepsilon>0\), \(x_\varepsilon=q_\beta-\varepsilon\) is a continuity point and satisfies \(F(x_\varepsilon)<\beta\).  Equation (16) and \(\eta_n\to0\) imply \(\Pr\{c_n>x_\varepsilon\}\to1\).  Equations (5), (17), and weak convergence then yield
\[
\liminf_{n\to\infty}\Pr\{T_n\leq c_n\}\geq F(x_\varepsilon).
\]
Letting \(\varepsilon\downarrow0\), and using the absence of an atom at the nonzero point \(q_\beta\), gives
\[
\liminf_{n\to\infty}\Pr\{T_n\leq c_n\}\geq F(q_\beta)\geq\beta.
\tag{18}
\]

It remains to handle \(q_\beta=0\), when
\[
F(0-)\leq\beta\leq F(0).
\tag{19}
\]
Fix a realization \(z=X_n\) satisfying \(H(z)=0\).  Since \(H(z)=\max\{D_L(z),-D_U(z)\}=0\), positive homogeneity and (6) give
\[
D_L(a_n)=\sqrt{m/n}\,D_L(z)\leq0,
\qquad
D_U(a_n)=\sqrt{m/n}\,D_U(z)\geq0.
\tag{20}
\]
For every \(g\in V_p\), these inequalities imply the global sign preservation
\[
\begin{aligned}
K_{a_n}(g)<0
&\Longrightarrow D_L(g)<D_L(a_n)\leq0
\ \text{and}\ D_U(g)>D_U(a_n)\geq0\\
&\Longrightarrow H(g)<0.
\end{aligned}
\tag{21}
\]
This step uses neither a selected projection cone nor a unique endpoint slope.

The event \(\{H(g)<0\}\) is a finite union of intersections of strict linear inequalities from the refined fan.  Approximate strict thresholds from below and apply the same finite inclusion--exclusion comparison used above.  It gives
\[
\Pr_p\{H(X_m)<0\}\leq F(0-)+C_F\delta_m(p).
\tag{22}
\]
Combining (12), (21), and (22), uniformly on \(\{\|X_n\|_2\leq M,H(X_n)=0\}\), yields
\[
\Pr^*\{S_{n,b}<0\}
\leq F(0-)+C_F\delta_m(p)
+\frac{m(m-1)}{2n}+\frac{m\sqrt d M}{2\sqrt n}+C_2e^{-c_2m}.
\tag{23}
\]

If \(F(0-)=\beta\), then \(\beta>0\) and \(H(G_p)<0\) with positive probability.  Hence at least one endpoint slope \(\ell\), nonzero on \(V_p\), occurs as a row of a matrix in \(\mathcal A_p\).  The scalar \(\ell^{\top}X_m\) has at most \((m+1)^d\) support points, so one of its jumps has size at least \((m+1)^{-d}\).  The law of \(\ell^{\top}G_p\) is continuous by nondegeneracy.  Comparing the two distribution functions immediately before and at that jump proves
\[
\delta_m(p)\geq\frac{1}{2(m+1)^d}.
\tag{24}
\]
Thus \(e^{-c_2m}=o(\delta_m(p))\).  Define
\[
R_n=\delta_{m_n}(p)+\frac{m_n}{\sqrt n}+\frac{m_n^2}{n}.
\]
The assumed rate condition is \(R_n=o(\eta_n)\).  If \(F(0-)=\beta\), equations (23)--(24) make its right-hand side \(\beta+o_{\Pr}(\eta_n)\); if \(F(0-)<\beta\), equations (14) and (23) instead leave a fixed positive gap below \(\beta\).  A lower \((\beta+\eta_n)\)-quantile of a finite distribution is nonnegative whenever the mass strictly below zero is less than \(\beta+\eta_n\).  Since \(X_n=O_{\Pr}(1)\), we conclude that
\[
\Pr\{H(X_n)=0,\ c_n<0\}\longrightarrow0.
\tag{25}
\]

We also have \(c_n=o_{\Pr}(1)\).  For every \(\varepsilon>0\), \(q_\beta=0\) implies \(F(-\varepsilon)<\beta\), so (16) gives \(c_n>-\varepsilon\) with probability tending to one.  Also \(F(\varepsilon)>\beta\): this is immediate if \(F(0)>\beta\); if \(F(0)=\beta\) and \(H\) takes a positive value, continuity, positive homogeneity, and full Gaussian support on \(V_p\) give a positive-probability open set on which \(0<H(G_p)\leq\varepsilon\); if \(H\) never takes a positive value, then \(F(\varepsilon)=1\).  Another application of (16) gives \(c_n\leq\varepsilon\) with probability tending to one.  Hence
\[
c_n=o_{\Pr}(1).
\tag{26}
\]

The zero set \(Z_H=\{g\in V_p:H(g)=0\}\) is a finite union of interiors of full-dimensional fan cones on which the selected linear functional is zero, together with subsets of finitely many zero and fan-facet hyperplanes.  Its boundary is therefore contained in finitely many hyperplanes with normals nonzero on \(V_p\).  Gaussian nondegeneracy gives \(\Pr\{G_p\in\partial Z_H\}=0\).  Thus \(Z_H\) is a continuity set for the weak convergence in (2), and
\[
\Pr\{H(X_n)=0\}\longrightarrow
\Pr\{H(G_p)=0\}=F(0)-F(0-).
\tag{27}
\]
For every \(\varepsilon>0\), equations (5), (25), and (26) ensure, up to an event of probability tending to zero, coverage of both disjoint events \(\{H(X_n)\leq-\varepsilon\}\) and \(\{H(X_n)=0\}\).  Equations (17) and (27) therefore imply
\[
\liminf_{n\to\infty}\Pr\{T_n\leq c_n\}
\geq F(-\varepsilon)+F(0)-F(0-).
\]
Letting \(\varepsilon\downarrow0\) makes the right-hand side \(F(0)\geq\beta\) by (19).  Together with (18), this proves
\[
\liminf_{n\to\infty}\Pr\{T_n\leq c_n\}\geq1-\alpha.
\tag{28}
\]
By (1), the event in (28) is precisely
\[
\left\{L(p)\geq\widehat L-c_n/\sqrt n,\quad
U(p)\leq\widehat U+c_n/\sqrt n\right\}.
\]
Because \(\eta_n\downarrow0\) and \(\alpha>0\), the level \(1-\alpha+\eta_n\) lies in \((0,1)\) for all sufficiently large \(n\), which is enough for the limiting assertion.

It remains to verify the stated thirty-two-type boundary certificate.  By \(\text{thm:complete-32-type-refinement-certificate}\), the table in \(\text{def:complete-32-type-table}\) is the complete observable--target quotient.  Every column of \(B^{32}\) has one entry \(1/2\) in each trial-arm block.  If a convex combination of columns equals
\[
p_\star=\tfrac12(\mathbf e_1+\mathbf e_7),
\]
every positively weighted column must be supported on rows \(1\) and \(7\), by nonnegativity.  Inspection of the complete table shows that precisely columns \(1\) and \(17\) have this support, that the two columns coincide, and that \(h^{32}_1=h^{32}_{17}=0\).  Hence \(p_\star\) is a vertex, its fiber is the segment joining types \(1\) and \(17\), and every point on that segment is optimal for both endpoint programs with value zero.  The multinomial covariance on the two positive cells gives
\[
G_{p_\star}=\zeta(\mathbf e_1-\mathbf e_7),
\qquad \zeta\sim\mathcal N(0,1/4).
\tag{29}
\]

We enumerate the projection on the two rays of \(V_{p_\star}=\operatorname{span}\{\mathbf e_1-\mathbf e_7\}\).  First suppose \(\zeta>0\), divide all candidate tangent coordinates by \(\zeta\), and consider mass leaving high-arm row \(7\).  Four outgoing high-arm cells, rows \(8,9,10,12\), are reachable while retaining low-arm row \(1\).  Let \(r\geq0\) be their total outgoing coordinate.  The remaining outgoing high-arm row \(11\), with coordinate \(t\geq0\), can only be reached by columns that also leave low-arm row \(1\) for row \(2\) or row \(5\).  Consequently the low-arm contribution to squared distance from \(\mathbf e_1-\mathbf e_7\) is at least
\[
(1+t)^2+t^2/2,
\tag{30}
\]
where the second term is the minimum squared norm obtained by splitting total mass \(t\) between two cells.  Splitting \(r\) among its four high-arm cells similarly shows that the high-arm contribution is at least
\[
(1-r-t)^2+r^2/4+t^2.
\tag{31}
\]
The sum of (30)--(31) is convex on \(r,t\geq0\), and its derivative with respect to \(t\) is \(2r+7t\).  Its minimum therefore has \(t=0\).  Minimization in \(r\) then gives \(r=4/5\), equally divided among the four reachable cells.  This lower bound is attained by the displayed table columns, proving
\[
\Pi_{T_{\mathfrak O^{32}}(p_\star)}
\{\zeta(\mathbf e_1-\mathbf e_7)\}
=\frac{\zeta}{5}(-4\mathbf e_7+\mathbf e_8+\mathbf e_9+\mathbf e_{10}+\mathbf e_{12}).
\tag{32}
\]
For \(\zeta<0\), the arm-swapped enumeration uses low-arm outgoing rows \(2,3,4,5\) and gives
\[
\Pi_{T_{\mathfrak O^{32}}(p_\star)}
\{\zeta(\mathbf e_1-\mathbf e_7)\}
=\frac{|\zeta|}{5}(-4\mathbf e_1+\mathbf e_2+\mathbf e_3+\mathbf e_4+\mathbf e_5).
\tag{33}
\]

On either ray, each of the four outgoing cells has a target-zero representative, and all target coordinates are nonnegative.  Hence
\[
D_L(G_{p_\star})=0.
\tag{34}
\]
Exactly two outgoing cells on each ray have target-one representatives whose other observed cell is the relevant base cell.  An incidence entry is \(1/2\), so producing outgoing observable derivative \(|\zeta|/5\) requires type-weight derivative \(2|\zeta|/5\).  Assigning each of the two eligible cells to its target-one representative attains total target derivative \(4|\zeta|/5\).  Conversely, nonnegative weights on nonbase columns cannot cancel an outgoing observable coordinate.  Thus only columns supported on a base cell and one of the four displayed outgoing cells can enter a tangent representation of (32) or (33), and table completeness shows that only the same two cells have a target-one such column.  This proves the matching upper bound and therefore
\[
D_U(G_{p_\star})=\frac45|\zeta|,
\qquad H(G_{p_\star})=0.
\tag{35}
\]

Because \(\mathfrak O^{32}\) is a polytope, a sufficiently small neighborhood of \(p_\star\) agrees with the corresponding neighborhood of \(p_\star+T_{\mathfrak O^{32}}(p_\star)\): in a finite inequality representation, active inequalities define the tangent cone and inactive inequalities retain positive slack.  If \(y=\Pi_{\mathfrak O^{32}}(p_\star+x)\), projection optimality against \(p_\star\) gives
\[
\|y-p_\star\|_2^2\leq x^{\top}(y-p_\star)
\leq\|x\|_2\|y-p_\star\|_2,
\]
and hence \(\|y-p_\star\|_2\leq\|x\|_2\).  Small inputs therefore project inside that common neighborhood, making (32)--(33) exact local formulas.  The target-zero and target-one allocations proving (34)--(35) remain feasible with weights linear in \(|\zeta|\), and the nonnegativity argument gives the matching local upper bounds.  The two rays form a finite fan; after vacuous whole-space equalities are discarded, its proper facet normals are nonzero on the one-dimensional space \(V_{p_\star}\).  Equation (29) is nondegenerate on that space.  Thus conditions (i)--(ii) hold even though both endpoint fibers have multiple optima, and (35) proves that the exact Gaussian directional law of \(H\) is the unit atom at zero.

Finally, (14) and \(m_n^2/n\to0\) imply
\[
r_n:=\delta_{m_n}(p_\star)+m_n/\sqrt n+m_n^2/n\longrightarrow0.
\]
The positive nonincreasing sequence
\[
\eta_n^\star=\sup_{k\geq n}\sqrt{r_k}+n^{-1}
\]
satisfies \(\eta_n^\star\downarrow0\) and
\[
0\leq\frac{r_n}{\eta_n^\star}\leq\sqrt{r_n}\longrightarrow0.
\]
Hence admissible slack schedules exist, and every decreasing schedule satisfying the theorem's displayed dominance condition yields the asserted pointwise conservative coverage.